%==============================================================
%  lecons/analyse/regularisation_zeta.tex — 1+2+3+… = -1/12 et effet Casimir
%==============================================================

\lecontitle{La somme \texorpdfstring{$1+2+3+\cdots=-\tfrac{1}{12}$}{1+2+3+... = -1/12}}{Analyse réelle — Régularisation zêta et effet Casimir}

%=== NOTATION LOCALE =========================================
% (aucune : \zeta, \hbar sont fournis par amsmath/amssymb)

%=============================================================
%  § 1 — DÉFINITION ET RÉSULTATS FONDAMENTAUX
%=============================================================
\secnum{navyblue}{1}{Définition et résultats fondamentaux}

\begin{tcolorbox}[thmbox]
  \textbf{Le paradoxe.}\enspace\index{régularisation zêta}
  La série $\sum_{n\geq 1} n = 1+2+3+\cdots$ \emph{diverge grossièrement}
  ($n\not\to 0$, sommes partielles $S_N=\tfrac{N(N+1)}{2}\to+\infty$). Pourtant
  l'égalité
  \[
    \boxed{\;1+2+3+4+\cdots \;\text{«\,$=$\,»}\; -\tfrac{1}{12}\;}
  \]
  apparaît en théorie des nombres (Ramanujan) et en physique (effet Casimir).
  Elle n'est \emph{pas} une somme au sens usuel : c'est la \emph{valeur
  régularisée} attachée à la série.

  \medskip
  \textbf{Définition (régularisation zêta).}\enspace\index{prolongement analytique}
  À la série $\sum_{n\geq 1} n^{a}$ on associe la fonction $s\mapsto\zeta(s)$,
  définie pour $\operatorname{Re}(s)>1$ par $\zeta(s)=\sum_{n\geq 1}n^{-s}$, puis
  \emph{prolongée analytiquement} à $\mathbb{C}\setminus\{1\}$. La \emph{valeur
  régularisée} de $\sum n^{a}$ est, par définition, $\zeta(-a)$. Ainsi
  $1+2+3+\cdots$ correspond à $a=1$, c'est-à-dire $\zeta(-1)$.

  \medskip
  \textbf{Théorème (valeurs régularisées).}
  \begin{itemize}
    \item \textit{Équation fonctionnelle.} On a l'identité entre fonctions
    méromorphes sur $\mathbb{C}$
    \[
      \zeta(s)=2^{s}\pi^{s-1}\sin\!\Bigl(\tfrac{\pi s}{2}\Bigr)\Gamma(1-s)\,\zeta(1-s),
    \]
    où, aux points exceptionnels ($s=0$ et $s$ entier $\geq 2$, où le membre de
    droite présente une forme indéterminée), l'identité s'entend par passage à
    la limite.
    \item \textit{Valeurs aux entiers négatifs.}
    \[
      \zeta(0)=-\tfrac12,\qquad
      \zeta(-1)=-\tfrac{1}{12},\qquad
      \zeta(-2k)=0\ (k\geq 1),\qquad
      \zeta(1-2k)=-\tfrac{B_{2k}}{2k},
    \]
    où $B_{2k}$ est le $2k$-ième nombre de Bernoulli.
    \item \textit{Cohérence.} La même valeur $-\tfrac1{12}$ s'obtient par d'autres
    procédés (sommation de Ramanujan, régularisation d'Abel via la fonction êta),
    ce qui en fait la valeur « canonique » associée à $1+2+3+\cdots$.
  \end{itemize}
\end{tcolorbox}

\begin{significationintuitive}
La série $1+2+3+\cdots$ vit dans une famille à un paramètre, $s\mapsto\sum
n^{-s}$, qui n'a de sens immédiat que pour $s>1$. Le prolongement analytique est
l'unique façon cohérente d'étendre cette fonction ; il « se souvient » de toute
la famille et attribue à $s=-1$ la valeur $-\tfrac1{12}$. Le nombre $-\tfrac1{12}$
n'est donc pas le résultat d'une addition, mais l'empreinte que laisse la série
divergente sur la fonction zêta — empreinte que la nature, via l'effet Casimir,
rend mesurable.
\end{significationintuitive}

\decsep{}

%=============================================================
%  § 2 — DÉMONSTRATION
%=============================================================
\secnum{navyblue}{2}{Démonstration}

\begin{ideepreuve}
On admet le prolongement analytique de $\zeta$ et son équation fonctionnelle
(établis par ailleurs). Le calcul de $\zeta(-1)$ en découle par évaluation
directe à l'aide de $\zeta(2)=\pi^2/6$. On confirme ensuite la valeur par un
chemin élémentaire : la fonction êta de Dirichlet $\eta(s)=(1-2^{1-s})\zeta(s)$,
combinée à la régularisation d'Abel de $1-2+3-4+\cdots=\tfrac14$ (établie dans
la leçon sur les séries divergentes), redonne $\zeta(-1)=-\tfrac1{12}$.
\end{ideepreuve}

\vspace{10pt}
\rubriq{Preuve}

\begin{mdframed}[style=proofbar]
  \textit{Étape 1 — Calcul de $\zeta(-1)$ par l'équation fonctionnelle.}
  On admet l'équation fonctionnelle. En $s=-1$,
  \[
    \zeta(-1)=2^{-1}\,\pi^{-2}\,\sin\!\Bigl(-\tfrac{\pi}{2}\Bigr)\,
      \Gamma(2)\,\zeta(2).
  \]
  Or $\sin(-\pi/2)=-1$, $\Gamma(2)=1!=1$ et $\zeta(2)=\pi^2/6$, donc
  \[
    \zeta(-1)=\frac{1}{2}\cdot\frac{1}{\pi^{2}}\cdot(-1)\cdot 1\cdot
      \frac{\pi^{2}}{6}
    =-\frac{1}{12}.
  \]

  \medskip
  \textit{Étape 2 — Confirmation par la fonction êta.}
  Pour $\operatorname{Re}(s)>1$, en séparant termes pairs et impairs,
  \[
    \eta(s)=\sum_{n\geq 1}\frac{(-1)^{n-1}}{n^{s}}
    =\Bigl(1-2^{1-s}\Bigr)\,\zeta(s),
  \]
  identité qui se prolonge analytiquement à $\mathbb{C}$. En $s=-1$,
  $1-2^{1-s}=1-2^{2}=-3$, d'où
  \[
    \eta(-1)=(-3)\,\zeta(-1),\qquad\text{soit}\qquad
    \zeta(-1)=-\frac{\eta(-1)}{3}.
  \]

  \medskip
  \textit{Étape 3 — Régularisation d'Abel de $\eta(-1)=1-2+3-4+\cdots$.}
  Considérons, pour $x\in[0,1[$, la série entière
  \[
    \sum_{n\geq 1}(-1)^{n-1}n\,x^{n-1}
    =\frac{\mathrm{d}}{\mathrm{d}x}\sum_{n\geq 1}(-1)^{n-1}x^{n}
    =\frac{\mathrm{d}}{\mathrm{d}x}\!\left(\frac{x}{1+x}\right)
    =\frac{1}{(1+x)^{2}}.
  \]
  Sa limite radiale en $x\to 1^{-}$ vaut $\tfrac14$ : la valeur régularisée
  (au sens d'Abel) de $1-2+3-4+\cdots$ est donc $\tfrac14$. On \emph{admet}
  ici que cette valeur d'Abel coïncide avec le prolongement analytique
  $\eta(-1)$ — c'est un théorème de type abélien, la série de Dirichlet de
  $\eta$ ne convergeant que pour $\operatorname{Re}(s)>0$ — de sorte que
  $\eta(-1)=\tfrac14$. Les étapes~2--4 constituent ainsi une confirmation
  cohérente de l'étape~1 plutôt qu'une seconde démonstration autonome.

  \medskip
  \textit{Étape 4 — Conclusion.}
  En reportant dans l'étape~2,
  \[
    \zeta(-1)=-\frac{\eta(-1)}{3}=-\frac{1/4}{3}=-\frac{1}{12}.
  \]
  Les deux chemins concordent : la valeur régularisée de $1+2+3+\cdots$ est
  $-\tfrac1{12}$.
  \hfill$\blacksquare$
\end{mdframed}

\begin{significationintuitive}
L'étape~1 est rigoureuse (modulo l'équation fonctionnelle admise) ; les
étapes~2--4 montrent qu'un procédé tout autre — Abel — tombe sur la même valeur.
C'est cette \emph{robustesse} qui distingue $-\tfrac1{12}$ d'un simple artefact :
plusieurs régularisations indépendantes s'accordent.
\end{significationintuitive}

\decsep{}

%=============================================================
%  § 3 — EXEMPLES, CONTRE-EXEMPLES ET EXERCICES
%=============================================================
\secnum{exgreen}{3}{Exemples, contre-exemples et exercices}

\rubriq{Exemples}

\begin{mdframed}[style=exbar]
  \textbf{1. L'effet Casimir.}\enspace\index{effet Casimir}%
  Deux plaques conductrices parallèles, distantes de $a$ dans le vide, sont
  \emph{attirées} l'une vers l'autre par les fluctuations quantiques du champ
  électromagnétique (Casimir, 1948). Chaque mode du champ porte une énergie de
  point zéro $\tfrac12\hbar\omega$ ; la somme de ces énergies diverge, mais sa
  \emph{régularisation} est finie. Dans le modèle unidimensionnel, les modes ont
  pour pulsations $\omega_n=\tfrac{\pi c}{a}n$, et l'énergie du vide entre les
  plaques s'écrit formellement
  \[
    E(a)=\frac{1}{2}\sum_{n\geq 1}\hbar\omega_n
        =\frac{\pi\hbar c}{2a}\sum_{n\geq 1} n
        \;\xrightarrow{\text{régularisation}}\;
        \frac{\pi\hbar c}{2a}\,\zeta(-1)=-\frac{\pi\hbar c}{24\,a}.
  \]
  La force $-\,\mathrm{d}E/\mathrm{d}a<0$ est attractive. C'est ici que
  $\zeta(-1)=-\tfrac1{12}$ devient physiquement mesurable.

  \smallskip\noindent
  \textbf{2. Le résultat tridimensionnel.}\enspace%
  Le calcul complet pour le champ électromagnétique entre deux plaques donne une
  énergie et une force par unité de surface
  \[
    \frac{E}{A}=-\frac{\pi^{2}\hbar c}{720\,a^{3}},\qquad
    \frac{F}{A}=-\frac{\pi^{2}\hbar c}{240\,a^{4}},
  \]
  confirmées expérimentalement (Lamoreaux, 1997, à quelques \% près).

  \smallskip\noindent
  \textbf{3. Sommation de Ramanujan.}\enspace\index{sommation de Ramanujan}%
  Dans sa lettre à Hardy (1913), Ramanujan attribue lui aussi $-\tfrac1{12}$ à
  $1+2+3+\cdots$, par un procédé fondé sur la formule d'Euler--Maclaurin : la
  « somme de Ramanujan » d'une série est le terme constant de son développement,
  et coïncide ici avec $\zeta(-1)$.
\end{mdframed}

\vspace{6pt}
\rubriq{Contre-exemples et pièges}

\begin{mdframed}[style=cexbar]
  \textbf{Ce n'est pas une somme.}\enspace%
  Écrire $1+2+3+\cdots=-\tfrac1{12}$ avec un signe d'égalité ordinaire est
  abusif : la série diverge, $S_N\to+\infty$. Le membre de droite est une valeur
  \emph{régularisée}, pas la limite des sommes partielles. Une somme de termes
  positifs — sommes partielles, moyennes ou limites — ne peut être négative.

  \smallskip\noindent
  \textbf{Manipulations illégales.}\enspace%
  La « preuve » populaire qui décale et soustrait des séries
  ($1-1+1-\cdots=\tfrac12$, puis $1-2+3-\cdots=\tfrac14$, puis soustraction)
  enchaîne des opérations interdites sur des séries divergentes : réarrangement,
  addition terme à terme. Elle \emph{suggère} le bon résultat mais ne le
  \emph{démontre} pas ; seule la régularisation (zêta, Abel, Ramanujan) lui donne
  un sens.

  \smallskip\noindent
  \textbf{Dépendance au procédé.}\enspace%
  Toutes les séries divergentes ne reçoivent pas une valeur, et la valeur peut
  dépendre du procédé pour des séries plus sauvages. Si $1+2+3+\cdots$ donne
  $-\tfrac1{12}$ par \emph{plusieurs} méthodes, c'est une propriété remarquable,
  pas une règle générale.

  \smallskip\noindent
  \textbf{La constante de coupure disparaît.}\enspace%
  En physique, l'énergie brute $\sum\tfrac12\hbar\omega_n$ est infinie ; seule la
  \emph{différence} avec le vide sans plaques est observable. La régularisation
  retire proprement la partie divergente (indépendante de $a$), et c'est le reste
  fini, en $\zeta(-1)$, qui porte la physique.
\end{mdframed}

\vspace{6pt}
\exolabel

\begin{mdframed}[style=exobar]
  \begin{enumerate}
    \item À partir de l'équation fonctionnelle, calculer $\zeta(-3)=\tfrac{1}{120}$
          (\textit{indication :} $\zeta(4)=\pi^4/90$). Retrouver ensuite
          $\zeta(0)=-\tfrac12$ : l'équation fonctionnelle y est \emph{indéterminée}
          (forme $0\cdot\infty$, car $\sin(\tfrac{\pi s}{2})\to 0$ tandis que
          $\zeta(1-s)\to\infty$ par le pôle en $1$), si bien qu'il faut lever
          l'indétermination par passage à la limite $s\to 0$ ; on pourra aussi,
          plus simplement, utiliser $\eta(0)=(1-2^{1-s})\zeta(s)\big|_{s=0}=-\zeta(0)$
          jointe à la valeur d'Abel $\eta(0)=1-1+1-\cdots=\tfrac12$.

    \item Démontrer $\eta(s)=(1-2^{1-s})\zeta(s)$ pour $\operatorname{Re}(s)>1$ en
          regroupant les termes pairs, et vérifier $\eta(-1)=-3\,\zeta(-1)$.

    \item Établir $\sum_{n\geq1}(-1)^{n-1}n\,x^{n-1}=(1+x)^{-2}$ pour $|x|<1$ et en
          déduire la valeur d'Abel $1-2+3-4+\cdots=\tfrac14$.

    \item Montrer que la « preuve » par décalage de séries viole la stabilité par
          réarrangement, en exhibant un réarrangement de $1-1+1-\cdots$ de moyenne
          de Cesàro différente de $\tfrac12$.

    \item (Modèle de Casimir 1D) Avec $\omega_n=\tfrac{\pi c}{a}n$, retrouver
          $E(a)=-\tfrac{\pi\hbar c}{24\,a}$ par régularisation zêta, puis calculer
          la force $F=-\,\mathrm{d}E/\mathrm{d}a$ et vérifier qu'elle est
          attractive.

    \item (Coupure exponentielle) Pour $\varepsilon>0$, calculer
          $\sum_{n\geq1} n\,e^{-\varepsilon n}$ et montrer que, lorsque
          $\varepsilon\to 0^{+}$,
          \[
            \sum_{n\geq1} n\,e^{-\varepsilon n}=\frac{1}{\varepsilon^{2}}
              -\frac{1}{12}+\mathcal{O}(\varepsilon^{2}).
          \]
          Interpréter $-\tfrac1{12}$ comme le terme \emph{fini} subsistant après
          retrait de la divergence $1/\varepsilon^2$.
  \end{enumerate}
\end{mdframed}

\leconfooter{S.\ Ramanujan (lettre à G.\,H.\ Hardy, 1913) ; H.\,B.\,G.\ Casimir (1948)}
